\documentclass[11pt]{article}
\usepackage[a4paper,margin=23mm]{geometry}
\usepackage{amsmath,amssymb,xcolor,hyperref,fancyhdr}
\usepackage{fontspec}
\setmainfont{Times New Roman}
\setmonofont{Times New Roman}
\usepackage[italic]{mathastext}
\setlength{\emergencystretch}{3em}
\setlength{\parskip}{7pt}\setlength{\parindent}{0pt}
\pagestyle{fancy}\fancyhf{}\setlength{\headheight}{14pt}
\lhead{FlowMLLab}\rhead{Reading and experiment guide}\cfoot{\thepage}
\begin{document}
\noindent\textbf{Week 7.1 | Reading the evidence}\par\medskip\hrule\medskip
\section*{Identify the model you actually run}
The default trained classroom model is a coordinate MLP, not a branch/trunk
DeepONet. An MLP maps a finite vector to outputs. A DeepONet combines a branch
encoding an input function sampled at fixed sensors with a trunk encoding query
coordinates. The optional topology demonstration is separate and may be skipped.

A classical interpolation baseline may win; reporting that is a correct result.
Locate errors by physical region as well as globally: shock layer, near-wall
region and wake. Quote the region mask and normalization. Similar training and
validation errors suggest limited fit capacity or budget, but do not identify a
unique cause. Neither a favorable error nor a topology diagram proves operator
generalization to a new family of input functions.
\paragraph{Check your understanding.}
Identify the input information, the reference, the metric denominator and the
evaluation split in one worked notebook comparison. Predict which conclusion
would change if any one of those four changed, then check your prediction.
\end{document}
